% Original applied explanation; see Chapter04/SOURCES.md.
\subsubsection{ক্লাবের নিবন্ধনে id, name ও value আলাদা কেন?}
ক্লাবের ফর্মে শিক্ষার্থীর নামের ঘরের \texttt{id} লেবেলের \texttt{for}-এর সঙ্গে মিলে কোন ঘরটির পরিচয় দেওয়া হচ্ছে তা বোঝায়। \texttt{name="student"} জমা দেওয়ার সময় তথ্যটির নাম এবং শিক্ষার্থী যা লিখেছেন তা তার বর্তমান মান। \texttt{placeholder} একটি ইঙ্গিত; এটি স্থায়ী লেবেলের বিকল্প নয়।

একটি বিভাগ বাছাইয়ের radio বোতামগুলোর \texttt{name} একই রাখলে তারা একই দল হয়, কিন্তু \texttt{value} আলাদা হলে নির্বাচিত বিভাগের পরিচয় যায়। একাধিক আগ্রহের জন্য checkbox উপযোগী। জমা দেওয়ার পর সার্ভারকে তথ্য যাচাই ও সংরক্ষণ করতে হয়; শুধু HTML ফর্ম বানালেই নিবন্ধনের তথ্য স্থায়ীভাবে জমা হয় না।
